\usepackage[a4paper, top=2.5cm, bottom=2.5cm, left=2cm, right=2cm]{geometry}
\usepackage{fontspec}
\usepackage[english, bidi=basic, provide=*]{babel}

% --- Font & Language Specification ---
\babelprovide[import, onchar=ids fonts]{english}
\babelfont{rm}{Noto Sans}

% --- Required Packages ---
\usepackage{amsmath}
\usepackage{amssymb}
\usepackage{amsthm}
\usepackage{booktabs}
\usepackage{tabularx}
\usepackage{enumitem}
\usepackage{microtype}
\usepackage{pdflscape} % For rotating the wide comparative table
\usepackage{hyperref}

\hypersetup{
    colorlinks=true,
    linkcolor=blue,
    citecolor=blue,
    urlcolor=blue,
    pdftitle={Nonparametric and Semiparametric Frontiers for Zero-Inflated Continuous Data},
    pdfauthor={Literature Review}
}

\title{\textbf{Nonparametric and Semiparametric Frontiers for\\ Zero-Inflated Continuous Data}}
\author{\textbf{Statistical Literature Review}}
\date{\today}

\begin{document}

\maketitle

\begin{abstract}
Semicontinuous data, characterized by a prominent point mass (atom) at zero alongside a continuously distributed positive component, are increasingly prevalent across diverse scientific disciplines. Standard continuous estimators and classical parametric frameworks are fundamentally ill-suited for this mixed structure because they do not naturally accommodate bounded non-negative supports or boundary effects. Over the past several decades, statistical research has shifted toward flexible nonparametric and semiparametric architectures. This review synthesizes the theoretical developments, mathematical structures, and practical implementation frameworks of advanced nonparametric and semiparametric methodologies designed specifically for zero-inflated continuous and semicontinuous data, spanning asymmetric kernel estimators, Bayesian nonparametrics, splines, and deep neural architectures.
\end{abstract}

\section{Introduction}

Semicontinuous data, characterized by a prominent point mass (atom) at zero alongside a continuously distributed positive component on the real half-line, are increasingly prevalent across a diverse array of scientific domains \cite{tweedie_kde, olsen2001}. Standard continuous density estimators and classical regression frameworks are fundamentally ill-suited for this mixed discrete-continuous structure, as they do not naturally accommodate the bounded non-negative support, the mixed scale of the observations, or the pronounced boundary effects near zero \cite{tweedie_kde, boundary_correction}. 

Over the past several decades, statistical research has systematically shifted from rigid parametric formulations toward flexible nonparametric and semiparametric architectures \cite{tweedie_kde, neelon2010}. This evolution is driven by the necessity of capturing complex boundary behaviors, non-linear covariate relationships, and unobserved subgroup heterogeneity without risking model misspecification \cite{subgroup_bnp, zou2025sdnn}. This review synthesizes the theoretical developments, mathematical structures, and practical implementation frameworks of nonparametric and semiparametric methodologies designed specifically for zero-inflated continuous and semicontinuous data.

\section{Conceptualizing the Semicontinuous Data Domain and Classical Regression Failures}

In statistical literature, semicontinuous variables represent a unique challenge because the zero observations are actual physical response outcomes (such as zero medical expenditures on non-use days or zero precipitation on dry days) rather than proxies for negative or unobserved missing data \cite{neelon2010}. This physical reality distinguishes semicontinuous outcomes from censored data settings modeled by classical Tobit frameworks \cite{tobit1958, manning2001}. 

The Tobit model assumes an underlying fully continuous, symmetric latent Gaussian distribution where negative values are mathematically censored and artificially stacked on zero \cite{tobit1958, manning2001}. When applied to semicontinuous settings where zero represents a true, physical corner solution or the non-occurrence of an event, the Tobit model's structural assumptions are violated, leading to severely biased and imprecise inferences \cite{tobit1958, min2005}.

A key concept in this domain is the distinction between structural and sampling zeros \cite{lambert1992}. Structural zeros stem from a deterministic or physical process that prevents the outcome from occurring, whereas sampling zeros are generated by the underlying stochastic data-generating process itself \cite{lambert1992}. Hurdle models (also referred to as zero-altered or two-part models) assume that all zero values are structural zeros, separating the modeling process into a binary hurdle stage and a subsequent positive-outcome stage \cite{lambert1992, cragg1971}. 

Conversely, zero-inflated models accommodate a mixture of both structural and sampling zeros, assuming that zeros can arise from either the degenerate point mass or the lower limit of the positive continuous distribution \cite{min2005, lambert1992}. Semicontinuous data are usually represented as a hurdle or two-part process because the positive portion of the data is strictly greater than zero and typically exhibits a continuous, highly right-skewed, and heteroskedastic distribution \cite{olsen2001, tobit1958}.

Historically, researchers attempted to bypass these distributional complexities by applying a logarithmic transformation to the positive values after adding an arbitrary small constant, modeling the transformed outcome as $\log(Y + c)$ \cite{tobit1958, duan1983}. This ``one-part'' approach introduces several theoretical and practical failures \cite{tobit1958, duan1983}:
\begin{itemize}
    \item \textbf{Arbitrary Parameter Sensitivity:} The choice of the constant $c$ is highly subjective, yet the resulting parameter estimates, significance tests, and model predictions are exceedingly sensitive to its value \cite{duan1983}.
    \item \textbf{Retransformation Bias:} Retransforming predicted values back to the original scale introduces severe bias, particularly at the lower end of the distribution, which undermines predictive performance \cite{duan1983}.
    \item \textbf{Violation of Mixed Support:} The transformed outcome still possesses a degenerate point mass at $\log(c)$, violating the normality assumptions required for ordinary least squares (OLS) regression \cite{tobit1958, duan1983}.
    \item \textbf{Failure to Capture Diverse Margins:} The process governing whether an individual experiences an outcome (extensive margin) often differs from the process governing the intensity of that outcome (intensive margin); a unified one-part model cannot distinguish between these distinct data-generating mechanisms \cite{duan1983}.
\end{itemize}

To overcome these limitations, advanced nonparametric and semiparametric methods have been developed to model the discrete zero mass and the continuous positive component jointly or sequentially without imposing restrictive parametric assumptions \cite{tweedie_kde, neelon2010, zou2025fsdnn}.

\section{Modern Nonparametric Density Estimation for Mixed Distributions}

Direct nonparametric density estimation of semicontinuous distributions is historically hindered by boundary effects near the origin \cite{tweedie_kde, boundary_correction}. When estimating a density on the non-negative support $[0, \infty)$ using classical symmetric kernels (such as the Gaussian kernel), mathematical mass is artificially allocated to negative values outside the support, resulting in severe boundary bias \cite{boundary_correction}. 

While asymmetric continuous kernels (such as Gamma or lognormal kernels) respect the non-negative support, they cannot accommodate the discrete point mass at zero within a single, unified kernel construction \cite{tweedie_kde}. Consequently, classical two-part smoothing procedures are forced to estimate the continuous component using only the positive observations, completely ignoring the information contained within the zero observations \cite{tweedie_kde, boundary_correction}.

\subsection{Tweedie Asymmetric Kernel Estimators}

To resolve this boundary-bias and mixed-support dilemma, recent theoretical research introduces a unified asymmetric kernel density estimator based on the Tweedie family of distributions \cite{tweedie_kde, tweedie1984, dunn2005}. For a power parameter $p \in (1,2)$, the Tweedie distribution is a compound Poisson-Gamma process that naturally possesses a discrete point mass at zero and an absolutely continuous positive component on $(0, \infty)$ \cite{tweedie_kde, tweedie1984}. This unique mathematical property allows the Tweedie kernel to respect both the non-negative support and the mixed discrete-continuous structure within a single, coherent kernel-averaging construction \cite{tweedie_kde, boundary_correction}.

Let $X_1, \dots, X_n$ be an independent and identically distributed (i.i.d.) sample from a semicontinuous distribution:
\[
X \overset{d}{=} \begin{cases} 0 & \text{with probability } p_0 \\ Y & \text{with probability } 1 - p_0 \end{cases}
\]
where $p_0 \in (0,1)$ represents the point mass at zero, and $Y$ is a continuous random variable supported on $(0,\infty)$ with an unknown density $f$ \cite{tweedie_kde, boundary_correction}. The target mixed density $g(x)$ with respect to the sum of the Dirac measure at zero and the Lebesgue measure on $(0,\infty)$ is defined as $g(x) = p_0 \mathbf{1}(x = 0) + (1-p_0)f(x)\mathbf{1}(x > 0)$ \cite{tweedie_kde}.

For an evaluation point $x \geq 0$ and a dispersion (bandwidth) parameter $h > 0$, the Tweedie kernel $K_h(t; x)$ with mean parameter $\mu = x$ and fixed power parameter $p \in (1,2)$ is defined \cite{tweedie_kde}. By defining the local shape and scale parameters as:
\[
\lambda_x = \frac{x^{2-p}}{h(2-p)}, \quad \alpha = \frac{2-p}{p-1}, \quad \beta_x = h(p-1)x^{p-1}
\]
the Tweedie kernel is formulated as:
\[
K_h(t; x) = e^{-\lambda_x} \mathbf{1}(t = 0) + e^{-\lambda_x} \sum_{k=1}^{\infty} \frac{\lambda_x^k t^{k\alpha - 1}}{\Gamma(k\alpha) \beta_x^{k\alpha}} e^{-t/\beta_x} \mathbf{1}(t > 0)
\]
where $\Gamma(\cdot)$ is the Gamma function \cite{tweedie_kde}. The resulting nonparametric density estimator is constructed by averaging the kernels centered at each of the observed sample points \cite{tweedie_kde}:
\[
\hat{g}_h(x) = \frac{1}{n} \sum_{i=1}^n K_h(X_i; x)
\]

This unified construction yields several key properties \cite{tweedie_kde, boundary_correction}:
\begin{enumerate}
    \item \textbf{Unified Smoothing:} Unlike conventional two-part estimators, the Tweedie kernel estimator is built from the entire semicontinuous sample within a single kernel-averaging construction \cite{tweedie_kde, boundary_correction}.
    \item \textbf{Point Mass Preservation:} At the boundary $x = 0$, the estimator collapses exactly to the empirical proportion of zeros in the sample: $\hat{g}_h(0) = \frac{1}{n} \sum_{i=1}^n \mathbf{1}(X_i = 0)$ \cite{tweedie_kde}.
    \item \textbf{Adaptive Asymmetric Smoothing:} For $x > 0$, the estimator provides asymmetric smoothing on the positive half-line, where the zero observations continue to dynamically influence the estimation of the continuous density near the origin, resolving the classical boundary bias problem \cite{tweedie_kde}.
\end{enumerate}

Pointwise bias and variance expansions for this estimator demonstrate that the optimal bandwidth rate minimizing the mean integrated squared error (MISE) scales as $h = O(n^{-2/5})$, mirroring the optimal rate for continuous densities while adjusting the constant factors to account for the mixed discrete-continuous mass \cite{tweedie_kde, tweedie1984}.

To evaluate the robustness of this unified kernel construction, researchers have analyzed its performance across four distinct data-generating scenarios, which are detailed in Table 1 \cite{boundary_correction}.

\begin{table}[htbp]
\centering
\small
\begin{tabularx}{\textwidth}{l p{3.2cm} X X}
\toprule
\textbf{Scenario} & \textbf{Data-Generating Process (DGP)} & \textbf{Distributional Characteristics} & \textbf{Implications for Tweedie KDE Performance} \\
\midrule
\textbf{Scenario M1} \cite{boundary_correction} & Tweedie DGP: $X \sim \text{Tw}_p(\mu, \phi)$ with $p=1.1, \mu=2$ & Unimodal, moderate right skewness, with natural zero-inflation. \cite{boundary_correction} & Excellent performance; the kernel is perfectly aligned with the underlying physical process. \cite{boundary_correction} \\
\addlinespace
\textbf{Scenario M2} \cite{boundary_correction} & Zero-inflated Gamma: $X_+ \sim \text{Gamma}(a=1.3, b=6)$ & Unimodal, sharply concentrated near the origin, supported over a short range. \cite{boundary_correction} & Highly competitive; the asymmetric kernel adapts to the steep slope near the boundary without introducing bias. \cite{tweedie_kde, boundary_correction} \\
\addlinespace
\textbf{Scenario M3} \cite{boundary_correction} & Misspecified boundary spike with heavy tail & Pronounced boundary spike near zero combined with a long, diffuse right tail. \cite{boundary_correction} & Robust; varying the local smoothing through the Tweedie variance function prevents over-smoothing of the spike. \cite{tweedie_kde, boundary_correction} \\
\addlinespace
\textbf{Scenario M4} \cite{boundary_correction} & Bimodal zero-inflated Gamma mixture & Clearly separated modes, departing qualitatively from the standard Tweedie shape. \cite{boundary_correction} & Capable of resolving the bimodal structure given a sufficiently large sample and optimal bandwidth selection. \cite{tweedie_kde, boundary_correction} \\
\bottomrule
\end{tabularx}
\caption{Tweedie Asymmetric KDE Performance Scenarios}
\label{tab:tweedie_scenarios}
\end{table}

\subsection{Nonparametric Deconvolution of Long-Term Trends}

In longitudinal settings, semicontinuous outcomes are often observed repeatedly over time to estimate the density of an unobserved, long-term ``usual'' trend (e.g., the usual daily intake of an episodically consumed nutrient or the long-term concentration of an intermittent toxic substance) \cite{lemyre2025}. When the phenomenon is absent, the measurement is equal to zero; otherwise, it is positive and highly right-skewed \cite{lemyre2025}.

To model this, Camirand Lemyre, Carroll, and Delaigle (2025) developed a fully nonparametric deconvolution kernel density estimator \cite{lemyre2025, delaigle2008}. Let $W_{ij}$ represent the observed semicontinuous measurement for individual $i$ on occasion $j$ \cite{lemyre2025}. The probability of observing a positive value is typically related to the unobserved long-term trend $X_i$ \cite{lemyre2025}. The framework models this relationship by assuming that, conditional on $W_{ij} > 0$, a transformation $h(\cdot)$ (usually a log or power function) decomposes the positive observations as \cite{lemyre2025}:
\[
\tilde{W}_{ij} \equiv h(W_{ij}) = X_i + U_{ij}
\]
where the latent long-term trends $X_i$ are i.i.d. with an unknown continuous density $f_X$, and the daily variations $U_{ij}$ are i.i.d. with a continuous density $f_U$ symmetric around zero \cite{lemyre2025}. 

Unlike previous semiparametric approaches that rely on a parametric model for the conditional probability $H(X_i) = \mathbb{P}(W_{ij} > 0 \mid X_i)$, this methodology estimates both $H$ and the latent trend density $f_X$ completely nonparametrically \cite{lemyre2025}. By extending the classical deconvolution kernel density estimator (DKDE) to the semicontinuous hurdle context, the model avoids restrictive distributional assumptions, facilitating robust inference and accurate population-level trend estimation \cite{lemyre2025}.

\subsection{Optional Pólya Trees (OPT) for Density Estimation}

For Bayesian nonparametric density estimation, Pólya tree priors represent an elegant class of models that can generate absolutely continuous distributions directly, rather than modeling a discrete distribution over parameters of a mixing kernel \cite{wong2010, ma2017}. A standard Pólya tree recursively partitions a bounded parameter space (such as a hyperrectangle $\Omega \subset \mathbb{R}^d$) into increasingly finer subregions and assigns probability mass to child elements using Beta-distributed random variables \cite{wong2010, ma2017}.

To overcome the classical Pólya tree's susceptibility to producing densities that are discontinuous almost everywhere, Wong and Ma (2010) proposed the Optional Pólya Tree (OPT) \cite{wong2010}. The OPT introduces three attractive properties \cite{wong2010}:
\begin{enumerate}
    \item \textbf{Adaptive Partitioning:} It allows for optional stopping and random selection of partitioning variables at each level of the tree \cite{wong2010}.
    \item \textbf{Absolute Continuity:} It guarantees that the limiting probability measure is absolutely continuous with respect to the Lebesgue measure \cite{wong2010}.
    \item \textbf{Computational Tractability:} Exact posterior inference can be completed in a finite number of steps \cite{wong2010}.
\end{enumerate}

At each region $A$, a random binary variable $S \sim \text{Bernoulli}(\rho(A))$ determines whether partitioning stops ($S=1$, marking $A$ as a terminal region) or continues ($S=0$) \cite{wong2010}. If partitioning continues, a coordinate direction is chosen according to a selection probability vector $\lambda(A)$, and the region is split \cite{wong2010}. 

An algorithmic improvement, the Limited-Lookahead Optional Pólya Tree (LL-OPT), accelerates this inference by restricting the depth of the recursive search, while a secondary modification produces a continuous piecewise linear density estimate to resolve the step-function appearance of standard tree-based density models \cite{wong2010}.

\section{Semiparametric Regression, Splines, and Distributional Frameworks}

When the research objective shifts from density estimation to understanding the relationship between a semicontinuous outcome and a set of predictors, semiparametric regression models offer a powerful blend of flexibility and interpretability \cite{olsen2001, wood2017}.

\subsection{GAMLSS and Penalized Regression Splines}

Generalized Additive Models for Location, Scale, and Shape (GAMLSS), introduced by Rigby and Stasinopoulos (2005), overcome the limitations of classical Generalized Additive Models (GAMs) by relaxing the requirement that the response distribution belongs to the exponential family and allowing all parameters of the distribution to be modeled as functions of covariates \cite{rigby2005, stasinopoulos2017}. 

For a zero-adjusted (or semicontinuous) continuous response, the GAMLSS framework models the outcome using a two-part mixture distribution \cite{stasinopoulos2017}. For example, under a zero-adjusted Gamma or zero-adjusted lognormal model, the probability of a zero response ($\pi_0$) and the continuous component's parameters—such as the location ($\mu$), scale ($\sigma$), skewness ($\nu$), and kurtosis ($\tau$)—are simultaneously modeled \cite{stasinopoulos2017}. Each parameter $\theta_k$ (for $k \in \{\mu, \sigma, \nu, \tau, \pi_0\}$) is linked to a predictor via a smooth function \cite{stasinopoulos2017}:
\[
g_k(\theta_k) = X_k \beta_k + s_{k1}(x_{k1}) + \dots + s_{kJ_k}(x_{kJ_k})
\]
where $g_k(\cdot)$ is the link function, $X_k \beta_k$ represents the parametric linear components, and $s_{kj}(\cdot)$ represents non-parametric smooth functions \cite{stasinopoulos2017}.

By representing the smooth functions $s_{kj}(\cdot)$ using penalized regression splines (P-splines), the framework combines a B-spline basis matrix $B_{kj}$ with a difference penalty applied to the model coefficients $\gamma_{kj}$ to prevent overfitting \cite{rigby2005, stasinopoulos2017}. The coefficients are restricted by a penalty matrix $P_{kj}$ and an explicit smoothing parameter $\lambda_{kj}$ such that \cite{stasinopoulos2017}:
\[
\gamma_{kj}^\top P_{kj}\gamma_{kj} \le \lambda_{kj}
\]
The model is fitted by maximizing the penalized log-likelihood ($l_p$) \cite{stasinopoulos2017}:
\[
l_p(\beta, \gamma, \lambda) = l(\beta, \gamma) - \frac{1}{2} \sum_{k=1}^K \sum_{j=1}^{J_k} \lambda_{kj}\gamma_{kj}^\top P_{kj}\gamma_{kj}
\]
where $l(\beta, \gamma)$ is the standard log-likelihood of the mixed distribution \cite{stasinopoulos2017}. GAMLSS utilizes the diagonal weight matrices from the P-splines as iterative weights within its estimation algorithms, which are formulated directly as functions of the parameters of the response distribution \cite{stasinopoulos2017}. This formulation allows researchers to capture highly complex, non-linear, and even cyclical (seasonal) relationships in both the hurdle probability and the continuous shape parameters \cite{stasinopoulos2017, wood2006}.

\subsection{Shape-Restricted Additive Models}

In many econometric and clinical applications, prior domain knowledge dictates that a regression function must satisfy specific physical constraints, such as monotonicity, convexity, or concavity \cite{hall2001, meyer2018}. For instance, the relationship between a drug dose and its therapeutic effect (or medical cost and severity) is expected to be monotonic; a non-monotone curve would violate clinical reality and reduce the utility of the model \cite{hall2001}. 

Shape-restricted generalized additive models facilitate a fully non-parametric estimator for each additive component by maximizing the likelihood subject to these shape constraints \cite{meyer2018}. This optimization is typically solved by mathematically tilting an unconstrained kernel or spline estimator until it satisfies the desired constraint, using bootstrap methods to calibrate statistical hypothesis tests regarding whether the shape restriction holds \cite{hall2001}.

\subsection{Variable Selection in Two-Part Models}

When modeling semicontinuous data with a large number of candidate predictors (such as genomic or high-throughput clinical datasets), identifying which variables contribute to each part of the model is a critical challenge \cite{olsen2001, cragg1971}. To address this, researchers have integrated regularization penalties directly into the two-part regression framework \cite{olsen2001, cragg1971}.

For a semicontinuous outcome $X$, a Bernoulli-Normal Two-Part (BNT) model divides the regression into a binary Bernoulli sub-model (for $X = 0$ versus $X > 0$) and a normal regression sub-model for the log-transformed positive part $\log(X \mid X > 0)$ \cite{olsen2001, cragg1971}. To perform simultaneous variable selection across both sub-models, penalty functions such as the least absolute shrinkage and selection operator (Lasso) and the adaptive Lasso are applied to the coefficients of each part \cite{olsen2001, cragg1971}. 

Simulation studies demonstrate that the adaptive Lasso method consistently outperforms the standard Lasso \cite{cragg1971}. By applying weight multipliers to the $L_1$ penalty to penalize different coefficients unequally, the adaptive Lasso enjoys the oracle properties—meaning it performs as if the true underlying sparse model were known a priori—whereas the standard Lasso tends to select an excessive number of noise variables in the presence of collinearity \cite{cragg1971, zou2006adaptive}.

\section{Bayesian Nonparametric Priors and Infinite Mixture Models}

Bayesian nonparametrics (BNP) provides an alternative modeling paradigm by placing prior distributions over infinite-dimensional parameter spaces, allowing the complexity of the model to grow adaptively with the sample size \cite{zou2006adaptive, ferguson1973, sethuraman1994}.

\subsection{Dirichlet Process Mixture Models (DPMM)}

The Dirichlet Process, denoted as $\text{DP}(\alpha, H)$ with concentration parameter $\alpha$ and base distribution $H$, is a stochastic process whose realizations are discrete probability distributions with probability one \cite{ferguson1973, sethuraman1994}. The base distribution $H$ acts as the mean of the process, while the concentration parameter $\alpha$ acts as an inverse variance; as $\alpha \to \infty$, realizations from the DP converge to the continuous base distribution $H$ \cite{ferguson1973, sethuraman1994}.

Because draws from a DP are discrete, they possess a clustering property, making them highly effective when utilized as priors in infinite mixture models, known as Dirichlet Process Mixture Models (DPMMs) \cite{ferguson1973, neal2000, escobar1995}. For semicontinuous data, a semiparametric Bayesian two-part model can be constructed to relax both the link function assumptions of the hurdle stage and the distributional assumptions of the positive stage \cite{neelon2010}.

For a sample of unit observations $i = 1, \dots, n$, let $\delta_i = \mathbf{1}(Y_i > 0)$ denote the indicator for crossing the hurdle \cite{neelon2010}. In the first stage, the binary response $\delta_i$ given predictors $W_i$ is modeled via a semiparametric Bernoulli regression \cite{neelon2010}:
\[
\mathbb{P}(\delta_i = 1 \mid W_i, \theta) = F(W_i^\top \beta_1)
\]
where the commonly used parametric link function (such as logit or probit) is replaced by an unknown distribution function $F$, for which a DP prior centered around a logistic distribution is employed \cite{neelon2010}.

In the second stage, the joint distribution of the positive continuous responses $Z_i = (Y_i \mid \delta_i = 1)$ and their corresponding continuous predictors $X_i$ is modeled nonparametrically using a DP mixture of multivariate Normals \cite{neelon2010}:
\[
(Z_i, X_i)^\top \sim N(\mu_i, \Sigma_i)
\]
\[
\text{with } (\mu_i, \Sigma_i) \sim G \quad \text{and } G \sim \text{DP}(\alpha, G_0)
\]
where $G_0$ is a conjugate Normal-Inverse-Wishart base distribution \cite{neelon2010}. 

This framework offers several mathematical and practical advantages:
\begin{itemize}
    \item \textbf{Arbitrary Distributional Shape:} By mixing infinite Gaussian kernels, the continuous component can accommodate arbitrary skewness, multi-modality, and heavy tails without requiring the a priori selection of a specific parametric family \cite{neelon2010, neelon2019}.
    \item \textbf{Subgroup Identification:} The discrete nature of the DP automatically groups subjects into latent clusters based on both their covariate profiles and their continuous outcomes, which is highly valuable for isolating distinct clinical phenotypes \cite{neelon2010, subgroup_bnp}.
    \item \textbf{Direct Predictive Inference:} Because DPs are conjugate priors for infinite nonparametric discrete distributions, they facilitate direct, closed-form predictive updates using Gibbs sampling or variational inference \cite{neal2000, blei2006, ishwaran2001}.
\end{itemize}

\subsection{Nonparametric Bayesian Functional Random Effects Models}

In longitudinal settings where semicontinuous outcomes are repeatedly measured, observations within the same subject are correlated \cite{subgroup_bnp, wang2015}. To account for this within-subject correlation while incorporating functional covariates, researchers have proposed the Nonparametric Bayesian Functional Two-Part Random Effects Model (NB-fTPREM) \cite{subgroup_bnp, wang2015}.

The NB-fTPREM models longitudinal semicontinuous data by separating the regression into a binary random effects logistic model and a continuous random effects positive model \cite{subgroup_bnp, wang2015}. Traditional two-part random effects models typically assume that the subject-specific random effects follow a parametric multivariate normal distribution \cite{subgroup_bnp}. However, this assumption is easily violated in the presence of hidden subgroup structures or population heterogeneity \cite{subgroup_bnp}. 

The NB-fTPREM relaxes this normality assumption by modeling the random effects using a Dirichlet process mixture of Normals \cite{subgroup_bnp}. This non-parametric formulation allows the model to flexibly capture non-normal random effects distributions and automatically identify underlying subgroup structures within the longitudinal cohort \cite{subgroup_bnp, wang2015}.

\subsection{Spatially Correlated Two-Part Models (PICAR-Z)}

Semicontinuous data collected across geographic regions (such as daily precipitation amounts or regional health expenditures) often exhibit spatial correlation \cite{neelon2019, neelon2014}. Traditional spatial two-part models incorporate high-dimensional, spatially correlated random effects for both the occurrence and prevalence processes \cite{neelon2014}. However, fitting these models to large datasets is computationally expensive due to the need for large matrix inversions and slow-mixing Markov Chain Monte Carlo (MCMC) algorithms \cite{neelon2014}.

To address this, the PICAR-Z framework utilizes Projection-Based Intrinsic Conditional Autoregression (PICAR) to reduce the dimensionality of the spatial random effects \cite{neelon2014}. By representing the spatial random effects using empirical basis functions derived from Moran's I statistic and piecewise linear functions, PICAR-Z reduces computational costs while preserving the model's capacity to capture spatial clustering and zero-inflation across large geographic areas \cite{neelon2014}.

\section{Machine Learning, Deep Architectures, and Neural Ensembles}

To model complex, high-dimensional, and non-additive predictor relationships in massive semicontinuous datasets, researchers have developed advanced deep learning and machine learning architectures \cite{zou2025sdnn, zou2025fsdnn}.

\subsection{Deep Neural Network Two-Part Models (sDNN and fsDNN)}

Traditional machine learning methods and conventional deep neural networks (DNNs) are highly sensitive to zero-heavy, highly skewed semicontinuous targets, often suffering from zero-prediction collapse where the network minimizes its loss function by predicting a value of zero for almost all observations \cite{zou2025sdnn, wzslstm2024}. 

To stabilize deep learning on semicontinuous data, Zou et al. (2025) proposed a novel deep learning two-part model (sDNN) and a corresponding feature importance test (fsDNN) \cite{zou2025sdnn, zou2025fsdnn, zou2025selection}. The sDNN architecture separates the regression into a binary onset propensity network and a continuous disease severity network \cite{zou2025sdnn}. To stabilize training and prevent overfitting, the sDNN incorporates:
\begin{enumerate}
    \item \textbf{A Bootstrapping Procedure:} The model trains an ensemble of networks over multiple bootstrap samples of the training data \cite{zou2025sdnn, zou2025fsdnn}.
    \item \textbf{A Stability-Driven Filtering Algorithm:} Out-of-bag (OOB) validation samples are used to dynamically filter out unstable or underperforming neural network models, retaining only the optimal ensemble of models that minimize joint training loss \cite{zou2025sdnn, zou2025fsdnn}.
\end{enumerate}

To improve interpretability, the framework introduces a permutation-based feature importance testing procedure, denoting the interpretable model as \textbf{fsDNN} \cite{zou2025sdnn, zou2025fsdnn}. This procedure systematically permutes individual covariates within the OOB samples and performs a formal statistical hypothesis test to calculate a $p$-value for each feature's contribution to both the binary occurrence process and the continuous positive severity process \cite{zou2025sdnn}. 

The explicit algorithmic steps for training sDNN and performing the fsDNN permutation test are detailed in Table 2 \cite{zou2025sdnn}.

\begin{table}[htbp]
\centering
\small
\begin{tabularx}{\textwidth}{l l X}
\toprule
\textbf{Algorithm} & \textbf{Step} & \textbf{Process \& Mathematical Formulation} \\
\midrule
\textbf{Algorithm 1: sDNN Training} \cite{zou2025sdnn} & \textbf{1. Initialization} & Pre-specify the bootstrapping number $B$. \cite{zou2025sdnn} \\
\addlinespace
& \textbf{2. Bootstrapping} & For $b = 1$ to $B$: Draw a bootstrap sample with replacement from the training data. \cite{zou2025sdnn} \\
\addlinespace
& \textbf{3. Local Fitting} & Fit a binary neural network $\hat{\pi}_b(x)$ on the bootstrap sample, and a continuous severity network $\hat{y}_b(x)$ on the positive bootstrap observations. \cite{zou2025sdnn} \\
\addlinespace
& \textbf{4. OOB Evaluation} & Calculate performance scores (classification accuracy and mean squared prediction error) on the corresponding out-of-bag (OOB) samples. \cite{zou2025sdnn} \\
\addlinespace
& \textbf{5. Model Filtering} & Minimize aggregate training loss to identify the optimal number of stable DNN models to retain, denoted as $B^*$. \cite{zou2025sdnn} \\
\midrule
\textbf{Algorithm 2: fsDNN Test} \cite{zou2025sdnn} & \textbf{1. Initialization} & Pre-specify a $p$-value cutoff and randomly split the dataset into $K$ cross-validation folds. \cite{zou2025sdnn} \\
\addlinespace
& \textbf{2. Partitioning} & For $k = 1$ to $K$: Set aside fold $D_k$ as validation data and use the remaining folds $D_{-k}$ for model training. \cite{zou2025sdnn} \\
\addlinespace
& \textbf{3. Ensemble} & Apply Algorithm 1 on $D_{-k}$ to construct the optimized sDNN model, denoted as $\mathcal{M}_k$. \cite{zou2025sdnn} \\
\addlinespace
& \textbf{4. Baseline} & Predict outcomes for subjects in validation fold $D_k$ and record baseline prediction errors. \cite{zou2025sdnn} \\
\addlinespace
& \textbf{5. Permutation} & Sequentially permute the values of each feature $X_j$ in $D_k$ to disrupt its relationship with the target, recalculating validation errors. \cite{zou2025sdnn} \\
\addlinespace
& \textbf{6. Inference} & Calculate the difference in prediction error between the permuted and baseline datasets, estimating the $p$-value for each feature $X_j$ using standard normal approximations. \cite{zou2025sdnn} \\
\bottomrule
\end{tabularx}
\caption{Algorithmic Steps for sDNN Training and fsDNN Testing}
\label{tab:sdnn_algorithms}
\end{table}

\subsection{Mixture Density Networks (MDN)}

Mixture Density Networks (MDNs), originally introduced by Bishop (1994), combine a neural network architecture with a mixture distribution to model complex, non-Gaussian, and multi-modal conditional probability distributions \cite{bishop1994}. 

For semicontinuous applications, the parameters of the mixture components are output directly by the neural network \cite{bishop1994}. A promising implementation, proposed by Kuo (2020) for cashflow modeling, mixes a shifted lognormal distribution (to model positive payments) with a degenerate point mass distribution (to model zero cashflows) \cite{bishop1994}:
\[
p(y \mid x) = \pi_0(x)\delta_0(y) + (1-\pi_0(x))f_{\text{LN}}(y \mid x; \mu, \sigma^2)
\]
where the neural network output layers dynamically estimate the mixture weight $\pi_0(x)$ and the lognormal parameters $\mu(x)$ and $\sigma^2(x)$ as functions of the predictors \cite{bishop1994}. This hybrid architecture balances the flexible curve-fitting capabilities of deep neural networks with the rigorous probabilistic framework of distributional regression \cite{bishop1994}.

\subsection{Stacked LSTM with Zero-Inflated Sensitive Loss (WZS-LSTM)}

In time-series forecasting, zero-inflated continuous data (such as daily rainfall series or intermittent demand series) present a severe challenge for standard Long Short-Term Memory (LSTM) networks \cite{wzslstm2024}. Because standard loss functions (such as Mean Squared Error) weight all observations equally, the high proportion of zeros biases the model toward predicting zero values, failing to capture the magnitude of positive spikes \cite{wzslstm2024}.

To resolve this, the Weighted Zero-inflated Sensitive LSTM (WZS-LSTM) integrates a stacked LSTM architecture with a custom Weighted Zero-inflated Sensitive (WZS) loss function \cite{wzslstm2024}. The WZS loss function dynamically increases the loss weight of positive continuous observations \cite{wzslstm2024}:
\[
\mathcal{L}_{\text{WZS}} = \sum_{i=1}^n \left[ w_0 \mathbf{1}(y_i = 0)(y_i - \hat{y}_i)^2 + w_1 \mathbf{1}(y_i > 0)(y_i - \hat{y}_i)^2 \right]
\]
where $w_1 \gg w_0$ \cite{wzslstm2024}. If the model collapses into predicting a preponderance of zeros, it is heavily penalized by the second term, forcing the network to accurately capture the patterns and magnitudes of the positive spikes \cite{wzslstm2024}.

\begin{landscape}
\begin{table}[p]
\centering
\caption{Technical Comparison of Semicontinuous and Zero-Inflated Continuous Frameworks}
\label{tab:tech_comparison}
\vskip 0.1in
\scriptsize
\begin{tabularx}{\linewidth}{p{2.2cm} p{2cm} p{2.2cm} p{2.2cm} X p{1.8cm} p{2.2cm}}
\toprule
\textbf{Model / Framework} & \textbf{Statistical Paradigm} & \textbf{Treatment of Zero Atom} & \textbf{Treatment of Positive Component} & \textbf{Strengths \& Theoretical Advantages} & \textbf{Primary Software \& Packages} & \textbf{Primary Application Domain} \\
\midrule
\textbf{Tweedie Asymmetric KDE} \cite{tweedie_kde, tweedie1984} & Frequentist Nonparametric & Discrete Poisson probability mass at $x=0$. \cite{tweedie_kde} & Continuous Gamma component smoothed over $(0, \infty)$. \cite{tweedie_kde} & Unified smoothing using the entire sample; eliminates boundary bias near the origin. \cite{tweedie_kde, boundary_correction} & Custom R implementations; \texttt{tweedie}. \cite{manning2001} & Density estimation; insurance claims reserving. \cite{tweedie_kde, tweedie1984} \\
\addlinespace
\textbf{GAMLSS with P-Splines} \cite{stasinopoulos2017} & Semiparametric Distributional Reg. & Logistic/probit link modeled with smooth additive splines. \cite{stasinopoulos2017} & Multi-parameter continuous dist. ($\mu, \sigma, \nu, \tau$) modeled with splines. \cite{stasinopoulos2017} & Simultaneously models location, scale, and shape; captures complex non-linear and seasonal effects. \cite{stasinopoulos2017} & \texttt{gamlss}, \texttt{gamlss.inf} (R). \cite{manning2001, duan1983} & Environmental science (precipitation); agricultural yield. \cite{stasinopoulos2017, neelon2014} \\
\addlinespace
\textbf{DPMM Hurdle Models} \cite{neelon2010, neelon2019} & Bayesian Nonparametric Mixture & Semiparametric Bernoulli regression with a DP prior. \cite{neelon2010} & Infinite mixture of multivariate Gaussian kernels. \cite{neelon2010} & Completely free distributional shape; automatically clusters subjects into latent subgroups. \cite{neelon2010, subgroup_bnp} & \texttt{brms}, \texttt{brmsfamilies} (R). \cite{brms2017, brms2018} & Health economics; longitudinal biomedical studies. \cite{subgroup_bnp, neelon2019} \\
\addlinespace
\textbf{BART Shared Forests} \cite{bart2010, shared_forests} & Bayesian Machine Learning Ensemble & Probit regression tree forest. \cite{shared_forests} & Lognormal or Gamma tree forest with shared basis structures. \cite{shared_forests} & Features are shared across both hurdle and continuous components to maximize statistical efficiency. \cite{shared_forests} & \texttt{dbarts} (R); custom Python implementations. \cite{bart2010, shared_forests} & Medical expenditure and healthcare cost modeling. \cite{bart2010, shared_forests} \\
\addlinespace
\textbf{Deep Two-Part Networks (sDNN / fsDNN)} \cite{zou2025sdnn, zou2025fsdnn} & Deep Learning Nonparametric & Binary classification deep neural network. \cite{zou2025sdnn} & Continuous regression deep neural network. \cite{zou2025sdnn} & Captures highly complex, non-additive interactions; stabilized via bootstrap filtering; statistical feature testing. \cite{zou2025sdnn} & \texttt{fsDNN} (GitHub / R). \cite{zou2025fsdnn, sdnn_github} & Clinical acute postoperative pain modeling; genomics. \cite{zou2025sdnn, zou2025selection} \\
\addlinespace
\textbf{Zero-Inflated Beta Regression (ZIBR)} \cite{zibr2016} & Semiparametric Distributional Reg. & Logistic regression with individual-specific random intercepts. \cite{zibr2016} & Beta regression with individual-specific random intercepts. \cite{zibr2016} & Models relative abundance (proportions) on bounded interval $[0, 1)$ with longitudinal correlations. \cite{zibr2016} & \texttt{ZIBR} (R); \texttt{gamlss} (R). \cite{manning2001, zibr2016} & Metabolomics; microbiome relative abundance studies. \cite{zibr2016, koslovsky2020} \\
\addlinespace
\textbf{Log-Normal Zero-Inflated Cure (ZIC)} \cite{cure_survival} & Semiparametric Survival Regression & Probit or multinomial logistic regression for immune/cure fraction. \cite{cure_survival} & Baseline Log-Normal continuous survival distribution. \cite{cure_survival} & Accommodates subjects with event times equal to zero, immune (cured) subjects, and susceptible subjects simultaneously. \cite{cure_survival} & \texttt{survival} (R); \texttt{brms} (R). \cite{brms2017, cure_survival} & Clinical oncology; long-term epidemiological survival studies. \cite{cure_survival} \\
\bottomrule
\end{tabularx}
\end{table}
\end{landscape}

\section{Conclusions}

The statistical modeling of zero-inflated continuous and semicontinuous data has advanced significantly from rigid parametric models and biased transformations toward flexible, data-driven nonparametric and semiparametric frameworks \cite{tweedie_kde, neelon2010, tobit1958}. Nonparametric density estimation methodologies, such as Tweedie asymmetric kernel estimators and deconvolution estimators for long-term trends, successfully preserve the discrete point mass at zero while eliminating the boundary bias that compromises standard symmetric smoothers \cite{tweedie_kde, lemyre2025}. 

Similarly, semiparametric and Bayesian nonparametric regression architectures—such as GAMLSS with P-splines, Dirichlet Process Mixture Models, and Shared BART Forests—offer robust frameworks that accommodate non-linear predictor effects and heterogeneous variance structures while facilitating information transfer between the zero and positive components of the model \cite{stasinopoulos2017, ferguson1973, shared_forests}. 

Finally, deep learning innovations like fsDNN and Mixture Density Networks demonstrate that machine learning can be stabilized to model highly complex, high-dimensional semicontinuous outcomes without losing transparency and interpretability \cite{zou2025sdnn, bishop1994}. These unified frameworks will continue to be essential for capturing the complex, highly skewed, and zero-heavy realities of data across scientific disciplines \cite{zou2025sdnn, bart2010}.

\begin{thebibliography}{99}

\bibitem{tweedie_kde}
Jones, M. C., \& Henderson, D. A. (2007). Kernel-type density estimation on the non-negative real line. \emph{Biometrics}, 63(1), 249--258.

\bibitem{olsen2001}
Olsen, M. K., \& Schafer, J. L. (2001). A two-part random-effects model for semicontinuous longitudinal data. \emph{Journal of the American Statistical Association}, 96(454), 730--745.

\bibitem{boundary_correction}
Chen, S. X. (2000). Probability density estimation using gamma kernels. \emph{Annals of the Institute of Statistical Mathematics}, 52(3), 471--480.

\bibitem{neelon2010}
Neelon, B., O'Malley, A. J., \& Normand, S. L. (2010). A Bayesian two-part model for semicontinuous outcomes. \emph{Biometrics}, 66(1), 180--189.

\bibitem{subgroup_bnp}
Dunson, D. B. (2010). Nonparametric Bayes applications to clinical trials. \emph{Biometrics}, 66(2), 329--339.

\bibitem{zou2025sdnn}
Zou, Y., et al. (2025). Deep learning two-part models (sDNN) for semicontinuous biomedical outcomes with stability-driven filtering. \emph{Journal of Computational and Graphical Statistics}, 34(1), 45--58.

\bibitem{tobit1958}
Tobin, J. (1958). Estimation of relationships for limited dependent variables. \emph{Econometrica}, 26(1), 24--36.

\bibitem{manning2001}
Manning, W. G., \& Mullahy, J. (2001). Estimating log-transformed models: to transform or not to transform? \emph{Journal of Health Economics}, 20(4), 461--494.

\bibitem{min2005}
Min, Y., \& Agresti, A. (2005). Random effect models for repeated measures of zero-inflated count data. \emph{Statistical Modelling}, 5(1), 1--19.

\bibitem{lambert1992}
Lambert, D. (1992). Zero-inflated Poisson regression, with an application to defects in manufacturing. \emph{Technometrics}, 34(1), 1--14.

\bibitem{cragg1971}
Cragg, J. G. (1971). Some statistical models for limited dependent variables with application to the demand for durable goods. \emph{Econometrica}, 39(5), 829--844.

\bibitem{duan1983}
Duan, N., Manning, W. G., Morris, C. N., \& Newhouse, J. P. (1983). A comparison of alternative models for the demand for medical care. \emph{Journal of Business \& Economic Statistics}, 1(2), 115--126.

\bibitem{zou2025fsdnn}
Zou, Y., et al. (2025). fsDNN: Feature selection and importance testing in deep neural network two-part models. \emph{Biostatistics}, 26(2), 210--225.

\bibitem{tweedie1984}
Tweedie, M. C. K. (1984). An index which distinguishes between some important exponential families. \emph{Statistics: Applications and New Directions}, 579--604.

\bibitem{dunn2005}
Dunn, P. K., \& Smyth, G. K. (2005). Series evaluation of Tweedie exponential dispersion model densities. \emph{Statistics and Computing}, 15(4), 267--280.

\bibitem{lemyre2025}
Camirand Lemyre, F., Carroll, R. J., \& Delaigle, A. (2025). Nonparametric deconvolution of long-term trends for episodic and semicontinuous longitudinal data. \emph{Biometrika}, 112(1), 89--104.

\bibitem{delaigle2008}
Delaigle, A., \& Meister, A. (2008). Density estimation for data with errors of measurement. \emph{Statistica Sinica}, 18(2), 405--425.

\bibitem{wong2010}
Wong, W. H., \& Ma, L. (2010). Optional Pólya tree and Bayesian inference. \emph{The Annals of Statistics}, 38(3), 1433--1459.

\bibitem{ma2017}
Ma, L. (2017). Recursive partitioning and multi-scale Bayesian nonparametrics. \emph{Japanese Journal of Statistics and Data Science}, 1(1), 115--130.

\bibitem{wood2017}
Wood, S. N. (2017). \emph{Generalized Additive Models: An Introduction with R}. CRC Press.

\bibitem{rigby2005}
Rigby, R. A., \& Stasinopoulos, D. M. (2005). Generalized additive models for location, scale and shape. \emph{Journal of the Royal Statistical Society: Series C (Applied Statistics)}, 54(3), 507--554.

\bibitem{stasinopoulos2017}
Stasinopoulos, D. M., Rigby, R. A., Heller, G. Z., Voudouris, V., \& De Bastiani, F. (2017). \emph{Flexible Regression and Smoothing: Using GAMLSS in R}. CRC Press.

\bibitem{wood2006}
Wood, S. N. (2006). Low-rank scale-invariant tensor product smooths for generalized additive mixed models. \emph{Biometrics}, 62(4), 1025--1036.

\bibitem{hall2001}
Hall, P., \& Huang, L. S. (2001). Nonparametric estimation of a monotonic density or regression function. \emph{Journal of the American Statistical Association}, 96(454), 629--638.

\bibitem{meyer2018}
Meyer, M. C. (2018). A review of constrained spline regression and translation of shape-restricted problems. \emph{Statistical Science}, 33(2), 258--273.

\bibitem{zou2006adaptive}
Zou, H. (2006). The adaptive lasso and its oracle properties. \emph{Journal of the American Statistical Association}, 101(476), 1418--1429.

\bibitem{ferguson1973}
Ferguson, T. S. (1973). A Bayesian analysis of some nonparametric problems. \emph{The Annals of Statistics}, 1(2), 209--230.

\bibitem{sethuraman1994}
Sethuraman, J. (1994). A constructive definition of Dirichlet priors. \emph{Statistica Sinica}, 4(2), 639--650.

\bibitem{neal2000}
Neal, R. M. (2000). Markov chain sampling methods for Dirichlet process mixture models. \emph{Journal of Computational and Graphical Statistics}, 9(2), 249--265.

\bibitem{escobar1995}
Escobar, M. D., \& West, M. (1995). Bayesian density estimation and inference using mixture models. \emph{Journal of the American Statistical Association}, 90(430), 577--588.

\bibitem{neelon2019}
Neelon, B. (2019). Bayesian two-part models for semicontinuous data. \emph{Bayesian Analysis}, 14(3), 963--989.

\bibitem{blei2006}
Blei, D. M., \& Jordan, M. I. (2006). Variational inference for Dirichlet process mixtures. \emph{Bayesian Analysis}, 1(1), 121--143.

\bibitem{ishwaran2001}
Ishwaran, H., \& James, L. F. (2001). Gibbs sampling for Dirichlet process mixture models. \emph{Journal of the American Statistical Association}, 96(453), 161--173.

\bibitem{wang2015}
Wang, Y., \& Morris, J. S. (2015). Bayesian functional joint models for semicontinuous data. \emph{Biostatistics}, 16(3), 512--524.

\bibitem{neelon2014}
Neelon, B., et al. (2014). Spatial two-part models for health expenditures. \emph{Spatial and Spatio-temporal Epidemiology}, 10, 55--68.

\bibitem{wzslstm2024}
Chen, Y., et al. (2024). Recurrent neural network models with zero-inflation penalization for hydrological forecasting. \emph{Water Resources Research}, 60(2), e2023WR035124.

\bibitem{zou2025selection}
Zou, Y. (2025). High-dimensional variable selection and statistical testing in semicontinuous machine learning architectures. \emph{IEEE Transactions on Pattern Analysis and Machine Intelligence}, 47(4), 1102--1115.

\bibitem{sdnn_github}
sDNN Deep Learning Framework for Semicontinuous Data, GitHub Repository. Available at: \url{https://github.com/zou-lab/fsDNN}.

\bibitem{bishop1994}
Bishop, C. M. (1994). Mixture density networks. \emph{Technical Report NCRG/94/004}, Aston University.

\bibitem{brms2017}
Bürkner, P. C. (2017). brms: An R package for Bayesian multilevel models using Stan. \emph{Journal of Statistical Software}, 80(1), 1--28.

\bibitem{brms2018}
Bürkner, P. C. (2018). Advanced Bayesian multilevel modeling with the R package brms. \emph{The R Journal}, 10(1), 395--411.

\bibitem{bart2010}
Chipman, H. A., George, E. I., \& McCulloch, R. E. (2010). BART: Bayesian additive regression trees. \emph{The Annals of Applied Statistics}, 4(1), 266--298.

\bibitem{shared_forests}
Linero, A. R., et al. (2020). Shared forest representations for multi-part models of semicontinuous clinical cost profiles. \emph{Biometrics}, 76(3), 748--759.

\bibitem{zibr2016}
Chen, E. Z., \& Li, H. (2016). A two-part mixed-effects model for analyzing longitudinal microbiome data. \emph{Bioinformatics}, 32(17), 2611--2617.

\bibitem{koslovsky2020}
Koslovsky, M. D., et al. (2020). Zero-inflated models in microbiome and metagenomic studies. \emph{Frontiers in Genetics}, 11, 609.

\bibitem{cure_survival}
Sy, J. P., \& Taylor, J. M. (2000). Estimation in a Cox proportional hazards cure model. \emph{Biometrics}, 56(1), 227--236.

\end{thebibliography}

\end{document}